\chapter{Diseño Hardware y Diagrama de Bloques Completo}
\label{ch:hardware_design}

\section{Arquitectura General y Esquema de Aislamientos Galvánicos}
\label{sec:arquitectura_general}

El concentrador de datos de nueva generación desarrollado en este Trabajo de Fin de Máster ha sido concebido como un nodo de procesamiento embebido de alta densidad, capaz de integrar en un único módulo compacto de carril DIN las funciones de telegestión de red por línea eléctrica (NB-PLC), metrología trifásica de precisión de Clase 0.2, comunicaciones ascendentes (WAN) y puertos de expansión inalámbricos.

La instalación del equipo en el entorno de un centro de transformación de media a baja tensión (MT/BT) exige un diseño de hardware altamente robusto frente a sobretensiones de Categoría IV y transitorios rápidos (\textit{bursts}/ESD). Por este motivo, el concepto de diseño pivota sobre una estricta \textbf{segmentación por dominios de masa y barreras de aislamiento galvánico}, tal y como se ilustra en la arquitectura global de la Figura \ref{fig:diagrama_bloques_global}.

\subsection{Estrategia de Aislamiento y Dominios de Masa}
\label{subsec:estrategia_aislamiento}

Para evitar bucles de masa, proteger los circuitos integrados de procesamiento de baja tensión ($3.3\text{ V}$ / $1.8\text{ V}$) y mitigar el acoplamiento de ruido en modo común procedentes de la red eléctrica, la placa base se divide en cuatro dominios galvánicos independientes:

\begin{enumerate}
    \item \textbf{Dominio de Alta Tensión y Potencia (HV-Domain):} Comprende la entrada de alimentación trifásica ($3\times230/400\text{ VAC} + \text{N}$), el filtro EMI primario y la etapa primaria de la fuente de alimentación conmutada de alta frecuencia (HVHF PSU). Este dominio soporta las tensiones de línea sin referencia a la masa digital del equipo.
    \item \textbf{Dominio de Medición y Metrología (MET-Domain):} Conecta directamente con los divisores de tensión resistivos de alta precisión y los transformadores de corriente (CT). Para garantizar la integridad de las medidas fasoriales e inmunitad ante sobretensiones, la interfaz de comunicación entre el chipset de metrología dedicado y el procesador principal se realiza mediante aisladores digitales magnéticos o capacitivos (SPI aislada), garantizando un aislamiento de al menos $2.5\text{ kV}_{\text{RMS}}$.
    \item \textbf{Dominio de Inyección y Frente Analógico PLC (AFE-Domain):} Incluye el driver amplificador de línea y los filtros analógicos de emisión/recepción. La inyección de la señal de comunicación por línea eléctrica sobre las fases de baja tensión requiere un acoplamiento capacitivo e inductivo sintonizado. El transformador de acoplamiento PLC proporciona la barrera de aislamiento galvánico de seguridad frente a los $230\text{ VAC}$ nominales, al tiempo que preserva una respuesta lineal en la banda de frecuencia CENELEC-A / FCC \cite{Hallak2013Coexistence}.
    \item \textbf{Dominio de Procesamiento y Comunicaciones SELV (Safe Extra Low Voltage):} Corresponde al núcleo digital alimentado a baja tensión continua ($+5\text{ VDC}$, $+3.3\text{ VDC}$). Incluye la unidad MPU con Linux embebido, la memoria RAM/eMMC, los transceptores Ethernet, las interfaces RS-485 aisladas, los optocopladores de E/S digitales y las ranuras de expansión para el módulo celular LTE y el módulo de radiofrecuencia (RF).
\end{enumerate}

\subsection{Interfaces de Campo y Protecciones Electromagnéticas}
\label{subsec:interfaces_protecciones}

Cada puerto físico que conecta con el exterior del chasis incorpora circuitos dedicados de protección frente a transitorios electromagnéticos:
\begin{itemize}
    \item \textbf{Puerto RS-485:} Implementa transceptores con aislamiento galvánico integrado y diodos de supresión de tensiones transitorias (TVS) en configuración diferencial, acompañados de resistencias PTC para protección ante cortocircuitos permanentes.
    \item \textbf{Puertos Ethernet (10/100 Mbps):} Utilizan conectores RJ45 con transformadores magnéticos integrados (\textit{magnetics}) que proporcionan un aislamiento galvánico de $1.5\text{ kV}_{\text{RMS}}$ conforme a la norma IEEE 802.3.
    \item \textbf{Entradas y Salidas Digitales (GPIOs de Campo):} Las entradas digitales de monitoreo (p. ej., detección de apertura de puerta del centro de transformación o alarmas de fusibles) y las salidas de relé están optocopladas mediante fototransistores de alta velocidad, aislando los transitorios de la planta del bus interno de la MPU.
\end{itemize}

\section{Subsistema de Procesamiento Central y Comunicaciones de Red}
\label{sec:procesamiento_central}

El núcleo de procesamiento del concentrador de datos requiere gestionar simultáneamente la pila de protocolos del nodo base (PRIME/G3-PLC), la pila de red IP/IPv6, el procesamiento en tiempo real de la metrología trifásica y la ejecución de servicios de ciberseguridad (TLS/IPsec).

\subsection{Unidad MPU y Controlador de Red (Linux Embebido)}
\label{subsec:mpu_linux}

Para responder a estos requerimientos se adopta una arquitectura basada en un System-on-Module (SoM) de alta eficiencia con microprocesador (MPU) de 32/64 bits ejecutando un sistema operativo Linux embebido industrial. 

Esta elección frente a microcontroladores de bajas prestaciones se justifica analíticamente por la latencia de procesamiento intrínseca de los modems NB-PLC basados en OFDM \cite{Bumiller2015Powerline}. Tal como demuestra Bumiller, las etapas de transformada rápida de Fourier (IFFT/FFT) y los algoritmos FEC (codificación convolucional concatenada con Reed-Solomon) introducen un retardo de cómputo de 3 a 5 símbolos OFDM por paquete. En estándares como IEEE 1901.2, este tiempo muerto o \textit{Response Interframe Spacing} (RIFS) alcanza hasta $695\,\mu\text{s}$ (equivalente a 10 símbolos OFDM en CENELEC-A) \cite{Bumiller2015Powerline}.

Para evitar el colapso por desbordamiento de memoria (\textit{buffer overflow}) en redes multisalto de alta densidad, la arquitectura de software del concentrador implementa un controlador a nivel de Kernel denominado \texttt{NetDevice} \cite{Bumiller2015Powerline}. Este driver abstrae la capa física mediante un protocolo de empaquetado a flujo (\textit{Packet-to-Stream Protocol}, P2SP) que multiplexa el tráfico de control, metrología y datos IP en flujos continuos, maximizando la eficiencia de los buffers de la MPU \cite{Bumiller2015Powerline}.

\subsection{Interfaces de Campo: Ethernet, RS-485 y E/S Digitales}
\label{subsec:interfaces_campo}

El procesador se comunica con el entorno del centro de transformación a través de interfaces aisladas:
\begin{itemize}
    \item \textbf{Doble Puerto Ethernet (10/100 Mbps):} Basado en controladores MAC integrados en la MPU y transceptores PHY externos comunicados vía MII/RMII. El aislamiento magnético de $1.5\text{ kV}_{\text{RMS}}$ previene interferencias galvánicas procedentes de switches de subestación.
    \item \textbf{Bus Industrial RS-485 Aislado:} Destinado a la conexión local de medidores modbus o unidades de supervisión de centro de transformación (TRM). Implementa transceptores con barrera galvánica integrada ($2.5\text{ kV}_{\text{RMS}}$) y terminación de línea conmutable mediante jumpers de hardware.
    \item \textbf{Entradas/Salidas Digitales Optocopladas:} Dos entradas digitales de detección ($0\text{--}24\text{ VDC}$) con aislamiento por optocoplador para monitorear alarmas físicas (p. ej., apertura de puerta del centro o disparo de fusibles) y dos salidas de relé de estado sólido (\textit{SSR}) para maniobra remota.
\end{itemize}

\section{Subsistema de Frente Analógico (AFE), PLC Trifásico y Comunicaciones Inalámbricas}
\label{sec:afe_plc_rf}

El subsistema de comunicaciones PLC constituye el elemento diferencial del concentrador de datos. Su diseño debe resolver el compromiso entre inyectar suficiente potencia en la red de baja tensión y garantizar la linealidad frente a las variaciones dinámicas de la impedancia de carga.

\subsection{Frente Analógico (AFE) y Driver de Inyección Trifásico}
\label{subsec:afe_driver}

El módulo analógico (AFE) se articula en torno a un módem módem transceptor integrado (p. ej., Microchip PL360 o Renesas CPX) gobernado por un DSP dedicado para la capa física PHY. Para inyectar la señal en las tres fases ($R, S, T$) sin sufrir los efectos severos de la desadaptación de impedancias descritos por Hallak y Bumiller \cite{Hallak2013Coexistence}, el hardware incorpora una etapa de amplificación lineal de alta corriente y acoplamiento trifásico direccionable.

En transmisiones OFDM multiportadora, la impedancia del circuito emisor en estado de transmisión ($TX$) debe ser extremadamente baja ($Z_1 \approx 0\,\Omega$) para garantizar la máxima transferencia de tensión hacia la red ($V_{TX}$) \cite{Hallak2013Coexistence}. Sin embargo, cuando operan múltiples transmisores no interoperables en la misma línea, esta baja impedancia en TX provoca la absorción y atenuación de la señal útil en fases adyacentes \cite{Hallak2013Coexistence}. Aunque la norma CENELEC EN 50065-7 fija restricciones de impedancia mínima en recepción/transmisión ($Z_e \ge 3\text{--}5\,\Omega$), su implementación física en etapas analógicas OFDM resulta inviable y costosa \cite{Hallak2013Coexistence}.

Para solventar esta limitación sin violar los límites de compatibilidad electromagnética, el hardware del concentrador utiliza un **esquema de acoplamiento capacitivo/inductivo trifásico con desacoplo de canal** \cite{Hallak2013Coexistence, Sendin2011Strategies}. Mediante una red de filtros pasa-banda sintonizados de alto orden y transformadores de aislamiento galvánico de alta frecuencia, el driver inyecta la señal de forma simétrica o seleccionable fase a fase, minimizando el acoplamiento parásito inter-fase y aislando la etapa digital de las tensiones de $230/400\text{ VAC}$.

\subsection{Ampliación a Banda Extendida (FCC) y Módulo Celular LTE}
\label{subsec:banda_fcc_lte}

Aunque la banda CENELEC-A (9--95 kHz) es el estándar de telegestión regulado en Europa, los ensayos empíricos de canal en entornos industriales demuestran que esta franja es altamente susceptible al ruido pulsante de motores y electrónica de potencia, registrando picos de interferencia de hasta $74\text{ dB}\mu\text{V}$ a $71\text{ kHz}$ \cite{Mishra2015Suitable}. Por el contrario, la banda FCC (10--490 kHz, centrada en $270\text{ kHz}$) presenta un suelo de ruido significativamente menor ($60\text{ dB}\mu\text{V}$), permitiendo alcanzar caudales de transmisión de hasta $34.28\text{ kbps}$ con modulaciones D8PSK y tasas de error de paquete ($\text{PER} < 1\%$) \cite{Mishra2015Suitable}. Por ello, el hardware del AFE se diseña con un ancho de banda flexible capaz de operar tanto en CENELEC-A como en la banda extendida FCC.

Para la red de transporte WAN hacia el sistema central (HES), la placa incluye un zócalo M.2 / mini-PCIe para la integración de un módem celular LTE Cat-1 / Cat-M1 / NB-IoT con soporte de tarjeta eSIM/SIM industrial y antena externa.

\subsection{Módulo RF Híbrido Inalámbrico (IEEE Std 1901.3-2025)}
\label{subsec:rf_hibrido}

Para solventar escenarios de atenuación extrema en la línea eléctrica (tales como la apertura de seccionadores o la presencia de filtros bloqueadores de alta impedancia), la arquitectura de hardware reserva una interfaz SPI/UART de alta velocidad para un módulo transceptor de radiofrecuencia (RF) sub-1 GHz.

Esta aproximación responde a las especificaciones del nuevo estándar **IEEE Std 1901.3-2025**, el cual establece el marco internacional para comunicaciones híbridas PLC/RF en redes inteligentes \cite{IEEE1901_3_2025}. La combinación síncrona de ambos medios en la capa MAC permite elevar la tasa de entrega de paquetes del $70\%$ (solo PLC) o $82\%$ (solo RF) a cifras superiores al **$90\%$** mediante enrutamiento híbrido dinámico \cite{Mishra2015Suitable}.

Conforme a la norma IEEE 1901.3-2025, el módulo RF opera en las bandas ISM ofreciendo tasas de velocidad de $25\text{ kbps}$ a $2\text{ Mbps}$, con modulaciones BPSK, QPSK y 16-QAM, y una sensibilidad de receptor de hasta $-109\text{ dBm}$ (modo MCS0) \cite{IEEE1901_3_2025}. El módulo actúa de forma transparente como estación o repetidor híbrido (\textit{Proxy Coordinator}, PCO), permitiendo que la MPU del concentrador alterne dinámicamente entre el canal PLC y el enlace inalámbrico según la métrica de calidad de enlace RSSI/SNR recibida \cite{IEEE1901_3_2025}.

